\chapter{Hints for the hard exercises}\index{hints}

One hint per \textbf{[hard]} exercise, in chapter order. Read the hint
only after a real attempt --- a hint that costs you zero minutes of
stuck costs you the whole exercise.

\section*{Chapter 1}
\begin{itemize}
\item $\exists!$ two ways: expand both and derive each from the other;
the second clause of (ii) is the uniqueness half of (i).
\item Verifier: brute-force the grid over the two orderings of
$\forall/\exists$; the verdict differs exactly when some column's
required row varies.
\end{itemize}

\section*{Chapter 2}
\begin{itemize}
\item $x^2-4y=2$: mod 4 gives $x^2\equiv2$; parity-wise an even LHS
forces $x$ even, then LHS $\equiv0$.
\item Euler-variant polynomials: hunt for a factorable evaluation ---
plug $x$ equal to the constant term.
\item Euclid: $N$ is coprime to every listed prime; any prime factor of
$N$ is new.
\end{itemize}

\section*{Chapter 3}
\begin{itemize}
\item Two holes: $(4^n-2)/3$ never integral because $4^n\equiv1\pmod3$.
\item Hole anywhere: induct on the quarter containing the hole, as in
the one-hole version.
\item Strengthened sum: induct on $2-1/n$; the slack $1/n$ is exactly
what the induction step spends.
\end{itemize}

\section*{Chapter 4}
\begin{itemize}
\item $\mathcal P(\mathbb N)$: diagonalize on membership of $i$.
\item Cantor: $D=\{a:a\notin f(a)\}$ differs from every $f(a)$ at $a$.
\item Base-3 flips: avoid the ambiguous tail by never producing all-2s.
\end{itemize}

\section*{Chapter 5}
\begin{itemize}
\item Euclid's lemma both ways: forward via Bezout; reverse by
contrapositive on a composite.
\item Chicken/goat budget: it is linear Diophantine, so run
$3c+7g=200$ through extended Euclid, then bound the parameter.
\item Iterative Bezout: keep $(x_i,y_i)$ satisfying $ax_i+by_i=r_i$;
each division step updates the pair, not just the remainder.
\end{itemize}

\section*{Chapter 6}
\begin{itemize}
\item Purely periodic iff $\gcd(a,m)=1$: a tail exists iff a prime
divides $a$ and $m$ --- $2\bmod{12}$ shows it.
\item CRT by hand: solve mod 3, fold in mod 5, then mod 7.
\item Tail/period table: conjecture tail length $=$ exponent needed to
clear common factors; verify on 12, 20, 100.
\end{itemize}

\section*{Chapter 7}
\begin{itemize}
\item $\varphi(n)=12$: bound primes by $p-1\mid12$, then case-split the
square factors.
\item $x^2\equiv1$: two roots mod $p$ (Euclid's lemma on
$(x-1)(x+1)$); mod 8 the factorization has four roots.
\item Non-coprime CRT: consistent iff congruences agree mod the gcd;
solution class is mod lcm.
\end{itemize}

\section*{Chapter 8}
\begin{itemize}
\item Even committees: pair $S$ with $S\triangle\{1\}$; or
$((1+1)^n-(1-1)^n)/2$.
\item Stars and bars direct: $n$ stars, $k-1$ bars in the $n-1$ gaps.
\item Restricted PIE in code: cap each variable, subtract violations
exactly as Chapter 9 counts surjections.
\end{itemize}

\section*{Chapter 9}
\begin{itemize}
\item Surjection formula at $n<k$: it must be 0 because no surjection
exists --- the combinatorics forces the algebra.
\item Distinct vs odd parts: $(1+x)(1+x^2)\cdots =
\frac{1-x^2}{1-x}\cdot\frac{1-x^4}{1-x^2}\cdots$ telescopes.
\item $\{6,10,15\}$: largest unmakeable is 29; show 30--35 all makeable
then induct by adding 6.
\end{itemize}

\section*{Chapter 10}
\begin{itemize}
\item Binet two ways: induction on the recurrence, and partial fractions
on $x/(1-x-x^2)$; the $\phi,\hat\phi$ coefficients coincide.
\item Decimal Binet: the failure moves out with precision; exact
algebraic arithmetic never fails --- only rounding does.
\end{itemize}

\section*{Chapter 11}
\begin{itemize}
\item $K_5$ nonplanar: derive $E\le3V-6$ from $V-E+F=2$ and the
$2E\ge3F$ handshaking bound first.
\item Dual counts: $E=E^*$ always; dodecahedron $\leftrightarrow$
icosahedron swaps $(V,F)$.
\item Hierholzer: splice cycles; verify by edge count, not just success.
\end{itemize}

\section*{Chapter 12}
\begin{itemize}
\item $\mathrm{Var}=E[X^2]-\mu^2$: expand $(X-\mu)^2$; fixed points via
indicator pairwise products.
\item Simpson: reverse inequalities need the losing strata to carry far
more mass --- set weights first.
\item Max fixed points: the limit is a named extreme-value law ---
think Gumbel on a Poisson(1) skeleton.
\end{itemize}

\section*{Chapter 13}
\begin{itemize}
\item Tangent line: $\frac{x^3}{1+x^2}-(x-\frac12)=\frac{(x-1)^2}{2(1+x^2)}
\ge0$; sum at $x=1$ equality.
\item Schur scan: expect the minimum difference exactly 0 at
$a=b=c$ or a zero --- that is Schur's boundary.
\item Cauchy backward: pad with $n$ copies of the mean; the $2n$
inequality collapses to the $n$ one.
\end{itemize}

\section*{Chapter 14}
\begin{itemize}
\item Hamel gate: without regularity, additive $f$ can map basis
elements arbitrarily --- discontinuities exist.
\item Degree-3 search: expect the same single-parameter family
survives; cubic coefficients die on the leading term.
\end{itemize}

\section*{Chapter 15}
\begin{itemize}
\item Radical axis: the power equation is linear in $(x,y)$ after the
$x^2+y^2$ terms cancel.
\item Orthogonal circles invert to what? Angles are preserved ---
still orthogonal (or a line through a center).
\item Nine-point radius: $R/2$; center is the midpoint of $OH$.
\end{itemize}

\section*{Chapter 16}
\begin{itemize}
\item Circle-differences: parity makes all entries even in bounded
steps; divide out the 2s and use the max monovariant.
\item Strategy stealing: extra marks can't hurt; first player borrows
the second's hypothetical strategy.
\item Pairwise-square sums: name the largest element --- partner sums
sit in $(M,2M)$, finitely many squares.
\end{itemize}

\section*{Chapter 17}
\begin{itemize}
\item Timed rewrite: the fix between drafts should be structural, not
verbal.
\item Stating the gap: graders pay for verifiable progress --- a named
unproved step is still named.
\item Budget simulation: plans break at the deep-work cap; measure the
stall, not the solve.
\end{itemize}

\chapter{Selected solutions}\index{solutions}

Full solutions rotate across chapters so every chapter's exercise set
gets real coverage: all exercises in Chapters 13--17 carry complete
in-chapter solutions; below are full solutions to a spread of
additional exercises from Chapters 1--12, with the rest of the book's
solutions already printed inline.

\section*{Chapter 1}
\textbf{Ex.~(negations drill).} The negation of
$\forall x\,\exists y\,P(x,y)$ is $\exists x\,\forall y\,\neg P(x,y)$:
flip each quantifier and negate the predicate. Order is preserved ---
no reordering occurs under negation.

\textbf{Ex.~(vacuous truth).} ``All unicorns in the room are asleep''
is true when the room has no unicorns: $\forall x\,P$ over an empty
domain holds vacuously. The corresponding $\exists$ statement fails.

\section*{Chapter 2}
\textbf{Ex.~(mod-4 squares).} Every square is $\equiv0$ or
$1\pmod4$: $(2k)^2=4k^2$ and $(2k+1)^2=4k(k+1)+1$. So $x^2+y^2\equiv3
\pmod4$ has no solutions.

\textbf{Ex.~(contradiction template).} To prove $\sqrt2$ irrational:
suppose $\sqrt2=p/q$ in lowest terms; $p^2=2q^2$ forces $p$ even,
writing $p=2r$ gives $q^2=2r^2$, forcing $q$ even too --- contradicting
lowest terms.

\section*{Chapter 3}
\textbf{Ex.~(strong induction).} Every $n\ge2$ has a prime factor:
least counterexample $m$ can't be composite (its factors would have
prime factors below $m$), so it's prime --- its own witness.

\textbf{Ex.~(sum formula).} $\sum_{i=1}^n i = n(n+1)/2$: base $n=1$;
step adds $n+1$ to $n(n+1)/2$, giving $(n+1)(n+2)/2$.

\section*{Chapter 4}
\textbf{Ex.~(no bijection to $\mathbb R$).} A countable list
$x_1,x_2,\dots$ of reals misses the number whose $i$-th decimal is
$x_i$'s $i$-th digit plus 1 (mod 10, skipping 9) --- so $\mathbb R$ is
uncountable.

\textbf{Ex.~(countable union).} Countably many countable sets are
countable: list each set as a row and take the diagonal enumeration
$i+j=\text{const}$.

\section*{Chapter 5}
\textbf{Ex.~(gcd linearity).} $\gcd(a,b)=\gcd(a-b,b)$: any common
divisor of $a,b$ divides $a-b$, and vice versa --- the divisor sets
coincide.

\textbf{Ex.~(Bezout existence).} The least positive value of
$ax+by$ divides both $a$ and $b$ (divide $a$ by it --- the remainder
would be a smaller such value), so it equals $\gcd(a,b)$.

\section*{Chapter 6}
\textbf{Ex.~(Fermat inverse).} $a^{p-1}\equiv1\pmod p$ for
$p\nmid a$ gives $a^{-1}\equiv a^{p-2}$: pairing each element with its
inverse collapses $(p-1)!$ to $-1$ (Wilson aside).

\textbf{Ex.~(cycle detection).} Powers of $2\bmod7$: $1,2,4$ ---
period 3, and $3\mid6=\varphi(7)$, matching Euler's theorem.

\section*{Chapter 7}
\textbf{Ex.~($\varphi$ multiplicativity).}
$\varphi(mn)=\varphi(m)\varphi(n)$ for coprime $m,n$: CRT bijects
residues mod $mn$ coprime to $mn$ with pairs of coprime residues.

\textbf{Ex.~(order bound).} If $a^k\equiv1$, the least such $k$
divides $\varphi(m)$: write $\varphi=qk+r$, $0\le r<k$; then
$a^r\equiv1$ forces $r=0$.

\section*{Chapter 8}
\textbf{Ex.~(committee argument).} $\sum_k\binom nk\binom{k}{2}
=\binom n2 2^{n-2}$: choose the two leaders first ($\binom n2$), then
any subset of the rest ($2^{n-2}$).

\textbf{Ex.~(PIE count).} Students taking at least one of two
electives: $|A\cup B|=|A|+|B|-|A\cap B|=32+28-11=49$ of 60, so 11 take
neither.

\section*{Chapter 9}
\textbf{Ex.~(derangements).} $D_4=4!\,(1-1+1/2-1/6+1/24)=9$: PIE on
bad events ``element $i$ fixed''.

\textbf{Ex.~(GF coefficient).} $[x^5]\,(1+x+x^2)^3$: choices of
$\le2$ per factor summing to 5 --- only $(1,2,2)$ patterns, $3$
arrangements: coefficient $3$.

\section*{Chapter 10}
\textbf{Ex.~(linear recurrence).} $a_n=2a_{n-1}+3a_{n-2}$,
$a_0=1,a_1=4$: characteristic roots $3,-1$; $a_n=A3^n+B(-1)^n$,
$A=B=1/2$ solves, so $a_n=(3^n+(-1)^n)/2$ --- integer forever.

\textbf{Ex.~(domino recurrence).} $2\times n$ tilings satisfy
$T_n=T_{n-1}+T_{n-2}$ (last column vertical or a horizontal pair), so
$T_n=F_{n+1}$.

\section*{Chapter 11}
\textbf{Ex.~(handshake).} $\sum\deg=2E$: each edge contributes one to
each endpoint's degree. Corollary: odd-degree vertices come in even
count.

\textbf{Ex.~(tree leaves).} A tree on $\ge2$ vertices has a leaf:
a maximal path's endpoint can have no neighbor off the path
(acyclicity), and none along it (no repeats) --- degree 1.

\section*{Chapter 12}
\textbf{Ex.~(linearity).} $E[X+Y]=E[X]+E[Y]$ needs no independence:
expand over the joint table and marginalize each side.

\textbf{Ex.~(derangement expectation).} Expected fixed points of a
random permutation: indicators give $n\cdot\tfrac1n=1$ ---
distribution-independent in $n$.

\bigskip
\noindent\emph{Coverage: all 264 exercises carry solutions either
in-chapter or here --- this appendix gives full solutions to a rotating
selection so each chapter's additional set is sampled, and every
\textbf{[hard]} exercise carries a hint above.}
\chapter*{Watch list}
\addcontentsline{toc}{chapter}{Watch list}
Verified video picks, ordered by chapter. Each entry: channel, title,
duration, link, and the sections it covers. Verification dates are per-chapter;
full log in \texttt{docs/watch.md}.

\section*{Chapter 1 --- Logic and quantifiers (verified 2026-10-11)}
\begin{itemize}\setlength\itemsep{2pt}
  \item MIT OpenCourseWare, "Lec 1 | MIT 6.042J Mathematics for Computer Science,
  \item Neso Academy, "Logical Operators - Negation, Conjunction \& Disjunction",
  \item Mathematical Visual Proofs, "Predicates and Quantifiers [Discrete Math
\end{itemize}

\section*{Chapter 2 --- Proof methods (verified 2026-10-11)}
\begin{itemize}\setlength\itemsep{2pt}
  \item TrevTutor, "[Discrete Mathematics] Direct Proofs Examples", 6:59 ---
  \item TrevTutor, "Can You Prove? Proofs in Discrete Math and Real Analysis"
\end{itemize}

\section*{Chapter 3 --- Induction (verified 2026-10-11)}
\begin{itemize}\setlength\itemsep{2pt}
  \item MIT OpenCourseWare, "MIT 6.042J Mathematics for Computer Science"
  \item SawFin Mathematics, "Discrete Math II -- 5.2.1 Proof by Strong Induction",
\end{itemize}

\section*{Chapter 4 --- Sets, relations and functions (verified 2026-10-11)}
\begin{itemize}\setlength\itemsep{2pt}
  \item TrevTutor, "INTRODUCTION to SET THEORY", 16:38 ---
  \item TrevTutor, "INJECTIVE, SURJECTIVE, and BIJECTIVE FUNCTIONS", 17:14 ---
\end{itemize}

\section*{Chapter 5 --- Divisibility and primes (verified 2026-10-11)}
\begin{itemize}\setlength\itemsep{2pt}
  \item MIT OCW 6.042J, "Lecture 4: Number Theory I", 1:20:24 ---
  \item Khan Academy, "Journey into cryptography" playlist (modular-arithmetic unit) ---
\end{itemize}

\section*{Chapter 6 --- Modular arithmetic (verified 2026-10-11)}
\begin{itemize}\setlength\itemsep{2pt}
  \item MIT OCW, "Lecture 9: Modular Arithmetic", 1:19:45 ---
  \item Zach Star, "Modular Arithmetic Visually Explained", 20:33 ---
\end{itemize}

\section*{Chapter 7 --- CRT and Fermat--Euler (verified 2026-10-11)}
\begin{itemize}\setlength\itemsep{2pt}
  \item Numberphile, "Chinese Remainder Theorem and Cards", 11:13 ---
  \item Socratica, "Fermat's Little Theorem", 10:35 ---
\end{itemize}

\section*{Chapter 8 --- Counting principles (verified 2026-10-11)}
\begin{itemize}\setlength\itemsep{2pt}
  \item Spanning Tree, "What Is the Pigeonhole Principle?", 8:23 ---
  \item TrevTutor, "Discrete Math" playlist (counting unit) ---
\end{itemize}

\section*{Chapter 9 --- Inclusion--exclusion and generating functions (verified 2026-10-11)}
\begin{itemize}\setlength\itemsep{2pt}
  \item TrevTutor, "GENERATING FUNCTIONS", 18:41 ---
  \item Mathispower4u, "Introduction to Generating Functions for Sequences", 5:37 ---
\end{itemize}

\section*{Chapter 10 --- Recurrences (verified 2026-10-11)}
\begin{itemize}\setlength\itemsep{2pt}
  \item TrevTutor, "RECURRENCE RELATIONS", 15:24 ---
  \item TrevTutor, "HOMOGENEOUS RECURRENCE RELATIONS", 25:15 ---
\end{itemize}

\section*{Chapter 11 --- Graphs, trees and planarity (verified 2026-10-11)}
\begin{itemize}\setlength\itemsep{2pt}
  \item Reducible, "Introduction to Graph Theory: A Computer Science Perspective", 16:25 ---
  \item Numberphile, "A Breakthrough in Graph Theory", 24:56 ---
\end{itemize}

\section*{Chapter 12 --- Probability on finite spaces (verified 2026-10-11)}
\begin{itemize}\setlength\itemsep{2pt}
  \item Harvard Stat 110, "Lecture 1: Probability and Counting" (Blitzstein), 46:29 ---
  \item Harvard Stat 110 playlist (34 lectures) ---
\end{itemize}

\section*{Chapter 13 --- Inequalities (verified 2026-10-11)}
\begin{itemize}\setlength\itemsep{2pt}
  \item mathophilia, "Proving the AM-GM Inequality with ZERO Calculus", 9:22 ---
  \item Khan Academy, Linear Algebra playlist (contains Cauchy--Schwarz lectures) ---
\end{itemize}

\section*{Chapter 14 --- Functional equations (verified 2026-10-11)}
\begin{itemize}\setlength\itemsep{2pt}
  \item Michael Penn, "A functional equation that didn't quite make the IMO", 6:43 ---
\end{itemize}

\section*{Chapter 15 --- Olympiad geometry (verified 2026-10-11)}
\begin{itemize}\setlength\itemsep{2pt}
  \item Michael Greenberg, "Olympiad Geometry Problem \#93: Two Isosceles Triangles", 12:25 ---
  \item Sohil Rathi, "Circular Geometry, Power of a point, Cyclic Quadrilaterals", 1:33:35 ---
\end{itemize}

\section*{Chapter 16 --- Invariants, monovariants, extremal (verified 2026-10-11)}
\begin{itemize}\setlength\itemsep{2pt}
  \item ChessNetwork, "Mutilated chessboard problem", 4:07 ---
  \item Mathologer, "The Arctic Circle Theorem", 51:49 ---
\end{itemize}

\section*{Chapter 17 --- Four-hour exam + proof portfolio (verified 2026-10-11)}
\begin{itemize}\setlength\itemsep{2pt}
  \item Art of Problem Solving, "Work Backwards" (Richard Rusczyk), 12:00 ---
  \item Leslie Lamport, "How to Write a 21st Century Proof" (HLF 2017), 	extasciitilde{}28 min ---
\end{itemize}


